\chapter{Tuberculosis (TB) Testing}

\section*{What Is This Test?}
Tuberculosis testing encompasses multiple diagnostic modalities for detection of \textit{Mycobacterium tuberculosis} complex infection. Testing includes: (1) Immunologic screening for latent TB infection (LTBI) using the tuberculin skin test (TST, also known as the Mantoux test or PPD) or interferon-gamma release assays (IGRAs, including QuantiFERON-TB Gold Plus and T-SPOT.TB), which measure the cell-mediated immune response to \textit{M. tuberculosis}-specific antigens (ESAT-6, CFP-10, TB7.7); (2) Acid-fast bacillus (AFB) smear microscopy of sputum or other clinical specimens for rapid presumptive identification of active pulmonary TB; (3) Nucleic acid amplification tests (NAATs), including the Xpert MTB/RIF Ultra (GeneXpert) assay, which detect \textit{M. tuberculosis} DNA and common rifampin resistance mutations (\textit{rpoB} gene) directly from sputum within 2 hours; (4) Mycobacterial culture (solid Lowenstein-Jensen medium or liquid MGIT 960) for definitive diagnosis and drug susceptibility testing (DST), requiring 1--8 weeks; and (5) Drug susceptibility testing (phenotypic and genotypic) for first-line drugs (rifampin, isoniazid, ethambutol, pyrazinamide) and second-line agents.

\section*{Alternative Names}
Tuberculosis Screening; TB Test; PPD; Mantoux Test; TST; IGRA; QuantiFERON-TB Gold; T-SPOT.TB; AFB Smear; AFB Culture; MTB PCR; GeneXpert MTB/RIF; Mycobacterium tuberculosis NAAT

\section*{Principle}
TST: 0.1 mL of purified protein derivative (PPD, 5 tuberculin units) is injected intradermally on the volar forearm, producing a 6--10 mm wheal. Induration (not erythema) is measured transversely at 48--72 hours. IGRA: Whole blood is incubated with \textit{M. tuberculosis}-specific antigens (ESAT-6, CFP-10), and interferon-gamma production by sensitized T-cells is measured by ELISA (QuantiFERON-TB Gold Plus) or enzyme-linked immunospot/ELISPOT (T-SPOT.TB). AFB smear: Ziehl-Neelsen (carbol fuchsin) or auramine-rhodamine fluorescent staining; results reported semi-quantitatively (scanty, 1+ to 4+). NAAT/Xpert MTB/RIF Ultra: Cartridge-based real-time hemi-nested PCR targeting the \textit{rpoB} gene with simultaneous detection of rifampin resistance-conferring mutations; limit of detection 15.6 CFU/mL. Mycobacterial culture: Processed (NALC-NaOH decontamination and concentration) specimens are inoculated into liquid medium (MGIT 960, BACTEC) and solid medium (Lowenstein-Jensen); growth detected by fluorescence (oxygen consumption in MGIT) or visual inspection. Drug susceptibility testing: Phenotypic (proportion method on solid medium, or MGIT 960 for liquid medium DST) or genotypic (line probe assays: GenoType MTBDRplus and MTBDRsl detecting mutations in \textit{rpoB}, \textit{katG}, \textit{inhA}, \textit{gyrA}, \textit{rrs}, \textit{eis} genes) for first- and second-line drugs.

\section*{History}
Robert Koch identified \textit{Mycobacterium tuberculosis} as the causative agent of tuberculosis in 1882, developed the Ziehl-Neelsen staining method, and produced tuberculin (old tuberculin). The tuberculin skin test was introduced by Clemens von Pirquet in 1907 and refined by Charles Mantoux in 1908. Florence Seibert developed purified protein derivative (PPD) in the 1930s. The WHO-recommended DOTS (Directly Observed Therapy, Short-course) strategy was launched in 1993. Interferon-gamma release assays (QuantiFERON-TB, then QuantiFERON-TB Gold, and T-SPOT.TB) were developed in the 2000s to improve specificity over TST by using antigens absent from BCG and most nontuberculous mycobacteria. The Xpert MTB/RIF assay (Cepheid) was endorsed by WHO in 2010 as a revolutionary rapid molecular diagnostic; the Ultra version with improved sensitivity was recommended in 2017. WHO's End TB Strategy (2015) targets a 90\% reduction in TB incidence by 2035. Multidrug-resistant TB (MDR-TB) and extensively drug-resistant TB (XDR-TB) remain critical global health threats, with approximately 450,000 incident MDR/RR-TB cases annually.

\section*{Why This Test Is Ordered}
\begin{itemize}
  \item Screening for latent TB infection in high-risk populations: close contacts of individuals with active pulmonary TB, persons born in or recently immigrated from countries with high TB prevalence (>20/100,000), residents and employees of high-risk congregate settings (correctional facilities, homeless shelters, nursing homes), healthcare workers, and individuals initiating immunosuppressive therapy (TNF-alpha inhibitors, organ transplant, prolonged corticosteroids)
  \item Diagnosis of active pulmonary TB in patients with persistent cough (>2--3 weeks), hemoptysis, fever, night sweats, weight loss, and abnormal chest radiograph (apical infiltrates, cavitary lesions, miliary pattern, pleural effusion)
  \item Rapid detection of pulmonary TB and rifampin resistance in patients with suspected drug-resistant TB (prior treatment, contact with MDR-TB, origin from high MDR-TB burden country)
  \item Diagnosis of extrapulmonary TB (lymphadenitis, pleural, CNS/meningitis, skeletal/Pott disease, genitourinary, peritoneal, pericardial) by AFB smear, NAAT, and culture of appropriate specimens
  \item Monitoring of anti-TB treatment response with serial sputum AFB smears and cultures at 2 months (end of intensive phase) and at treatment completion
  \item Required screening prior to initiation of biologic immunosuppressive agents (infliximab, adalimumab, etanercept) that are associated with reactivation of latent TB
\end{itemize}

\section*{Primary vs Secondary Test}
For latent TB infection, IGRA or TST serves as the primary screening test. A positive result must be followed by clinical evaluation and chest radiography to exclude active TB disease. For active TB, AFB smear and NAAT are primary rapid diagnostic tests, with culture serving as the definitive (reference) standard. NAAT (Xpert MTB/RIF) is the WHO-recommended initial diagnostic test in adults and children with suspected pulmonary TB and/or rifampin resistance.

\section*{How This Test Is Performed}
TST: 5 TU of PPD (0.1 mL) is injected intradermally on the volar surface of the forearm using a 27-gauge needle with the bevel upward; a raised wheal of 6--10 mm confirms intradermal (not subcutaneous) injection. The patient returns at 48--72 hours for measurement of induration transversely across the forearm. IGRA: Blood is collected into specialized tubes (nil, TB antigen, mitogen for QuantiFERON; or heparin tube for T-SPOT) and must be incubated at 37\textdegree{}C within 16 hours (QuantiFERON) or processed within 8 hours (T-SPOT). AFB sputum: Three consecutive early-morning sputum specimens collected 8--24 hours apart (at least one first-morning specimen) are recommended for initial diagnosis. Sputum induction with hypertonic saline may be used if spontaneous expectoration is inadequate. For children or patients unable to produce sputum, gastric aspirate (collected in the early morning before the patient rises) or bronchoalveolar lavage (BAL) may be obtained. For extrapulmonary TB, appropriate specimens include CSF, pleural fluid, lymph node aspirate/biopsy, urine, tissue biopsy, or bone marrow.

\section*{What Preparation Is Needed}
TST: No special preparation. The patient must be able to return at 48--72 hours for reading. IGRA: No fasting required. Blood must be processed in a timely manner. AFB sputum: Rinse mouth with water before collection; early morning specimens are preferred. Patients should be informed that three separate sputum collections are required for optimal sensitivity. Food and drink do not affect the specimen but may be withheld before sputum induction due to aspiration risk.

\section*{Contraindications and Precautions}
TST: Contraindicated in persons with documented prior positive TST (risk of ulceration and necrosis), prior severe allergic reaction to PPD, or extensive burns/rash on both forearms. A prior BCG vaccination can cause false-positive TST results (especially if vaccinated after infancy or with multiple BCG doses), but IGRA is not affected by BCG. TST is subject to the booster effect (two-step testing recommended in serial testing programs to distinguish boosting from true conversion). IGRA: False-negative results can occur in immunosuppressed patients (HIV with CD4 <200 cells/$\mu$L, anti-TNF therapy, chronic renal failure, high-dose corticosteroids). An indeterminate IGRA (low mitogen response) suggests anergy and does not exclude TB infection. AFB smear: A negative smear does NOT exclude TB, as smear sensitivity is only 50--60\%; 10,000 AFB/mL are required for a positive smear. NAAT should not be used to monitor treatment response as it detects DNA from both viable and nonviable organisms.

\section*{Risks}
TST: Local pruritus, erythema, and induration are expected. Rarely: vesiculation, ulceration, or necrosis at the injection site, particularly in persons with strong reactions. IGRA: Standard venipuncture risks (pain, bruising, hematoma, phlebotomy site infection). AFB sputum collection: No significant risks. The greatest risk associated with TB testing is a false-negative result leading to failure to diagnose active TB with resulting transmission to contacts.

\section*{What Happens After the Test}
TST must be read at 48--72 hours; results read after 72 hours may underestimate induration and are not valid. IGRA results available in 24--48 hours. AFB smear results available within 24 hours (STAT) or 1--3 days (routine). Xpert MTB/RIF Ultra results available in 2 hours. Mycobacterial culture results: preliminary positive (MGIT) at 1--3 weeks; final negative at 6--8 weeks. All positive or suspected TB results must be reported immediately to the ordering provider and the local public health department (TB is a nationally notifiable disease). Patients with suspected active pulmonary TB should be placed in airborne infection isolation (negative-pressure room, N95 respirator) until TB is excluded or the patient is no longer considered infectious (3 negative AFB sputum smears, clinical improvement, and >2 weeks of effective therapy).

\section*{Understanding Results}

\subsection*{Below Normal}
\begin{itemize}
  \item \textbf{TST negative (<5 mm induration):} No evidence of TB infection in most individuals. However: False-negative results occur in immunosuppressed patients; a negative TST does not exclude TB infection, especially in the setting of HIV, organ transplant, or anti-TNF therapy. Anergy testing is no longer recommended. A negative TST in a recent contact (within 8 weeks of exposure) may represent the pre-conversion window period
  \item \textbf{IGRA negative:} Below the manufacturer cutoff. QuantiFERON-TB Gold Plus: TB Antigen minus Nil <0.35 IU/mL. T-SPOT.TB: $\leq$4 spots (with Nil $\leq$10 spots). Indeterminate results (Nil >8.0 IU/mL or Mitogen minus Nil <0.5 IU/mL) require repeat testing
  \item \textbf{AFB smear negative:} Does not exclude TB. 50--60\% of culture-positive TB is smear-negative. Paucibacillary disease, early infection, extrapulmonary TB, and pediatric TB are commonly smear-negative
  \item \textbf{NAAT (Xpert) negative:} Does not completely exclude TB, particularly in paucibacillary disease. Sensitivity of Xpert Ultra in smear-negative culture-positive TB is 63--83\%. If clinical suspicion is high, culture should be obtained
  \item \textbf{Culture negative at 6--8 weeks:} No growth of \textit{M. tuberculosis} complex. Strongly suggests absence of active TB (NPV >99\% with adequate specimen)
\end{itemize}

\subsection*{Above Normal}
\begin{itemize}
  \item \textbf{TST positive:} Induration $\geq$5 mm (immunosuppressed, close contacts, fibrotic changes on CXR), $\geq$10 mm (recent immigrants, residents/employees of high-risk facilities, children <4 years, mycobacteriology lab workers), or $\geq$15 mm (persons with no known risk factors). Indicates TB infection but does NOT distinguish latent from active disease. Requires CXR and clinical evaluation
  \item \textbf{IGRA positive:} Indicates TB infection. May be reported as positive, negative, or indeterminate. Like TST, does not distinguish latent from active disease
  \item \textbf{AFB smear positive (scanty, 1+, 2+, 3+, 4+):} Reports the presence of acid-fast organisms visualized microscopically (ZN stain: red/pink rods against blue background; auramine-rhodamine: fluorescent yellow-orange rods against dark background). 1+ = 1--9 AFB/100 fields; 2+ = 1--9/10 fields; 3+ = 1--9/field; 4+ = >9/field. Smear positivity correlates with bacillary load and infectiousness. Smear-positive patients are 5--10 times more infectious than smear-negative patients
  \item \textbf{NAAT (Xpert MTB/RIF Ultra) positive:} Reports MTB detected (with semi-quantitation: trace, very low, low, medium, high) and rifampin resistance detected/not detected/indeterminate. Rifampin resistance-positive results should be confirmed by phenotypic DST and line probe assay
  \item \textbf{Culture positive for \textit{M. tuberculosis} complex:} Definitive diagnosis of active TB. Species identification (MALDI-TOF or DNA probe) differentiates \textit{M. tuberculosis} from nontuberculous mycobacteria (NTM). Drug susceptibility testing should follow
\end{itemize}

\section*{Normal Results}
TST: \textbf{No induration} (<5 mm; cutoff depending on risk group). IGRA: \textbf{Negative}. AFB smear: \textbf{No acid-fast bacilli seen}. NAAT: \textbf{MTB not detected}. Mycobacterial culture: \textbf{No \textit{Mycobacterium tuberculosis} complex isolated at 6--8 weeks}.

\section*{Follow-up Tests}
\begin{itemize}
  \item \textbf{Chest radiograph (PA and lateral):} Required for all patients with positive TST or IGRA to evaluate for active TB disease; CXR findings (apical infiltrates, cavitary lesions, hilar adenopathy, pleural effusion, miliary pattern) guide further diagnostic evaluation
  \item \textbf{Xpert MTB/RIF Ultra:} WHO-recommended initial test for individuals with suspected pulmonary TB; automatically reflexed when AFB smear is positive, or performed directly on the specimen
  \item \textbf{Mycobacterial culture with DST:} Phenotypic first-line DST (rifampin, isoniazid, ethambutol, pyrazinamide) for all initial culture-positive isolates; second-line DST if MDR-TB is suspected or confirmed
  \item \textbf{Line probe assay (GenoType MTBDRplus/MTBDRsl):} Molecular detection of mutations conferring resistance to rifampin, isoniazid, fluoroquinolones, and injectable second-line agents from culture isolates or direct smear-positive specimens
  \item \textbf{HIV testing:} Recommended for all patients with diagnosed TB infection (latent or active), as co-infection significantly impacts management
  \item \textbf{Liver function tests (ALT, AST, bilirubin):} Baseline and monthly monitoring during treatment with isoniazid, rifampin, and pyrazinamide due to hepatotoxicity risk
\end{itemize}

\section*{References}
\begin{enumerate}
  \item Lewinsohn DM, Leonard MK, LoBue PA, et al. Official ATS/IDSA/CDC clinical practice guidelines: diagnosis of tuberculosis in adults and children. Clin Infect Dis. 2017;64(2):e1-e33.
  \item World Health Organization. WHO consolidated guidelines on tuberculosis. Module 3: Diagnosis -- rapid diagnostics for tuberculosis detection. WHO; 2021.
  \item Nahid P, Dorman SE, Alipanah N, et al. Official ATS/CDC/IDSA clinical practice guidelines: treatment of drug-susceptible tuberculosis. Clin Infect Dis. 2016;63(7):e147-e207.
  \item Pai M, Behr MA, Dowdy D, et al. Tuberculosis. Nat Rev Dis Primers. 2016;2:16076.
  \item Steingart KR, Schiller I, Horne DJ, et al. Xpert MTB/RIF assay for pulmonary tuberculosis and rifampicin resistance in adults. Cochrane Database Syst Rev. 2014;(1):CD009593.
\end{enumerate}

\clearpage
\section*{Regulatory Status and Authorization Date}
Regulatory status applies to the specific assay, instrument, manufacturer, and intended use rather than to the test name alone. No single approval date can be assigned to this test category. For a commercial assay, verify the current product in the relevant regulator's database: the U.S. Food and Drug Administration (FDA) clearance, approval, or authorization; India's Central Drugs Standard Control Organisation (CDSCO) licence or permission; the European Union In Vitro Diagnostic Regulation (EU IVDR) CE certificate issued through the applicable conformity-assessment route; or the United Kingdom Medicines and Healthcare products Regulatory Agency (MHRA) registration and UKCA/CE status. Laboratory-developed tests may not have an individual FDA, CDSCO, EU IVDR, or MHRA approval date.
\clearpage
